\section{The Causal Simulation Harness}
\label{ch:simulation}

Every mission-level claim in this report rests on a single experimental
instrument: a time-stepped world simulator that exercises the complete
pipeline, from raw link observations through constrained edge pruning,
graph neural network scoring, and task assignment, down to task outcomes
and trust updates. The defining design requirement of this harness is that
mission success must \emph{causally} depend on the pruned topology, not
merely correlate with it. A purely correlational design, in which outcomes
are drawn independently of which links survived pruning, would make the
comparison between pruners meaningless: any pruning strategy would look as
good as any other, and a simulator could not tell the difference between
a method that preserves mission-critical connectivity and one that severs
it. We therefore built the harness around an explicit causal chain in
which the kept graph determines how reliably a task result can be
delivered, and delivery reliability feeds directly into whether the task
counts as a success. This chapter describes that harness: the per-cycle
phases, the hidden ground truth against which decisions are graded, the
delivery model that closes the causal loop, the fault and staleness
machinery, and the statistical protocol attached to every comparison.

\subsection{Per-Cycle Phases}

The simulation advances in discrete cycles. Within each cycle, five
phases execute in a fixed causal order, so that every quantity a phase
consumes was produced by a phase that ran before it in the same cycle or
in an earlier one. Figure~\ref{fig:sim-phases} summarises the flow.

\begin{figure}[htbp]
  \centering
  \begin{tikzpicture}[
      font=\small,
      phase/.style={rectangle, rounded corners=2pt, draw, thick,
        align=center, minimum width=4.1cm, minimum height=1.05cm,
        fill=white},
      fb/.style={rectangle, rounded corners=2pt, draw, thick, dashed,
        align=center, minimum width=4.1cm, minimum height=0.9cm},
      arr/.style={-{Stealth[length=2.5mm]}, thick},
      node distance=0.55cm and 1.6cm
    ]
    \node[phase] (p1) {\textbf{Phase 1: World step}\\ rovers move, links observed\\ (packet loss $5\,\%$), freshness\\ decays, batteries drain};
    \node[phase, below=of p1] (p2) {\textbf{Phase 2: Candidate graph}\\ build dynamic graph, compute\\ fresh per-edge features};
    \node[phase, below=of p2] (p3) {\textbf{Phase 3: Pruning}\\ QUBO pruner or baseline selects\\ the working graph};
    \node[phase, below=of p3] (p4) {\textbf{Phase 4: Scoring and}\\ \textbf{assignment}\\ GNN scores on the working graph,\\ tasks assigned, faults may act};
    \node[phase, below=of p4] (p5) {\textbf{Phase 5: Outcomes}\\ sample success from the hidden\\ reliability field};
    \node[fb, right=of p3] (fb1) {\textbf{Trust update}\\ feeds next cycle's\\ node and edge features};

    \draw[arr] (p1) -- (p2);
    \draw[arr] (p2) -- (p3);
    \draw[arr] (p3) -- (p4);
    \draw[arr] (p4) -- (p5);
    \draw[arr, dashed] (p5.east) -- ++(0.8,0) |- (fb1.south);
    \draw[arr, dashed] (fb1.north) |- ($(p1.east)+(0.8,0)$) -- (p1.east);
  \end{tikzpicture}
  \caption{The five per-cycle phases of the simulation harness, executed
    in causal order. The dashed loop closes the cycle: task outcomes
    update the trust model, and the updated trust enters the node and
    edge features of the following cycle.}
  \label{fig:sim-phases}
\end{figure}

\paragraph{Phase 1: world step.} Rovers first move, which changes pairwise
proximities and hence which links are physically available at all. Each
link within communication range is then observed subject to a packet loss
process with a default loss probability of $5\,\%$. A link that is not
observed in a cycle does not vanish instantly; instead, its freshness
decays with an exponential time constant of $3$ cycles, so that recent but
unconfirmed information is down-weighted rather than discarded. Batteries
drain during the same phase. This phase produces the raw, noisy,
partially stale view of the world that a real fog-tier fleet would
actually possess.

\paragraph{Phase 2: candidate graph construction.} From the current
observations the harness builds the dynamic candidate graph and attaches
fresh per-edge features: trust, link quality, battery level, proximity,
observation freshness, task criticality, and an estimate of processing
cost. These are exactly the features the constrained QUBO formulation of
the earlier chapters consumes; nothing is hidden from or added to the
pruner beyond what the observation model provides.

\paragraph{Phase 3: pruning.} A pruner selects the working graph from the
candidate graph. The production choice is the constrained QUBO pruner,
solved with the projected annealer. For comparison, the harness can
instead keep the full graph, apply a link-quality threshold
($\tau = 0.4$), or retain the top fraction ($0.5$) of edges per node.
All pruners see the same candidate graph and the same features, so any
difference in downstream mission performance is attributable to the
pruning decision itself.

\paragraph{Phase 4: scoring and assignment.} The trust-aware graph neural
network runs one message-passing round \emph{on the working graph only}
and scores every rover. Tasks, two per cycle by default, are then
assigned by capacity-aware greedy matching. Scripted fault events may be
active during this phase, degrading the true reliability of specific
rovers for a defined window.

\paragraph{Phase 5: outcomes and trust update.} Task outcomes are sampled
from the hidden ground-truth reliability field described below. The
observed successes and failures update the per-rover trust model, and the
updated trust values become features in the next cycle. This feedback
loop is essential: the harness evaluates not a static estimator but a
closed-loop system whose beliefs evolve as evidence accumulates.

\subsection{Hidden Ground Truth}

Each rover possesses a true task-execution reliability that the pipeline
never observes directly. The truth is structured in two deliberate ways.

First, it is \emph{spatially correlated}: nearby rovers share terrain and
interference conditions, so their reliabilities are drawn from a common
regional level. This is what gives the graph structure real content. If
reliabilities were independent, a neighbour's track record would carry no
information about a rover's own, and graph-based scoring would have no
legitimate advantage to discover. Under spatial correlation, message
passing over the working graph aggregates genuinely informative evidence.

Second, the truth may be \emph{nonstationary}. An optional per-cycle
random walk lets each rover's reliability drift over the mission. With
drift enabled, a trust value that was accurate ten cycles ago is stale
now, which is precisely the regime in which observation freshness,
exploration, and the probation treatment of re-added edges earn their
keep. All headline results in this report state explicitly whether drift
was active.

Two further properties of the ground truth deserve emphasis. First, it
is generated before any pipeline decision is made within a cycle, so no
modelling choice can leak into the standard against which it is graded.
Second, every run is fully determined by its seed: the same seed
reproduces the same reliability field, the same packet loss realisations,
and the same outcome draws, which is what makes the paired per-seed
statistics of Section~\ref{sec:sim-stats} meaningful. When two arms are
evaluated on the same set of seeds, they face literally identical worlds
and differ only in the decisions they take inside them.

\subsection{The Causal Delivery Model}

The causal heart of the harness is the rule that converts a pruning
decision into a mission outcome. A task counts as successful only if two
independent events both occur:

\begin{enumerate}
  \item \textbf{Execution.} The assigned rover executes the task
    successfully, a Bernoulli event whose probability is that rover's
    hidden reliability.
  \item \textbf{Delivery.} The result is delivered back to the root of
    the network through the \emph{kept} graph.
\end{enumerate}

Delivery probability is governed by a \emph{path factor}: the widest-path
(maximum bottleneck) reachability from the assigned rover to the root,
computed over the working graph with link quality as the bottleneck
metric. Intuitively, a result travels along the route whose weakest link
is strongest; if every remaining route crosses a poor link, delivery is
unlikely, no matter how reliable the executing rover was.

This model is what makes the harness honest about pruning. Because the
path factor is computed on the \emph{kept} graph, a pruner that removes
redundant links can strictly reduce the set of available routes and
thereby degrade delivery for exactly the missions that depend on
bottleneck-limited paths. The harness is therefore fully capable of
reporting that the QUBO pruner loses to a denser baseline on raw mission
success, and indeed it does so at larger fleet scales; we treat that
outcome as a finding to explain, not a defect to tune away. A simulator
whose outcomes ignored the pruned topology could never surface such a
result, which is why we regard the causal delivery model as the single
most important property of the evaluation.

\subsection{Default Configuration}

Table~\ref{tab:sim-params} collects the key harness parameters and their
defaults. These values define the canonical configuration; every
experiment that deviates from them says so.

\begin{table}[htbp]
  \centering
  \caption{Key parameters of the simulation harness and their default
    values.}
  \label{tab:sim-params}
  \begin{tabular}{@{}llp{5.6cm}@{}}
    \toprule
    Parameter & Default & Role \\
    \midrule
    Fleet size & 6 rovers & Number of rovers in the arena \\
    Tasks per cycle & 2 & Tasks offered for assignment each cycle \\
    Communication range & $4.0$\,m & Radius within which links can exist \\
    Packet loss & $0.05$ & Per-cycle link observation failure \\
    Freshness time constant & 3 cycles & Decay rate of unobserved links \\
    Exploration rate & $0.15$ & Fraction of assignments used to probe \\
    Load increment & $0.15$ & Compute load added per assigned task \\
    Load decay & $0.5$ & Per-cycle recovery of available compute \\
    Minimum trust & $0.25$ & Floor below which trust never falls \\
    Compute need & $0.3$ & Minimum free compute a task requires \\
    Mission-critical nodes & 3 & Nodes connectivity must always reach \\
    QUBO budget & 15 & Global budget in the pruning constraint \\
    \bottomrule
  \end{tabular}
\end{table}

\subsection{Fault Injection and Edge Lifecycle}

Controlled degradation is as important to the evaluation as nominal
operation. The harness supports \emph{scripted fault events}, each with a
start cycle and a duration. A typical event drops one rover's hidden
reliability for a fixed window, after which the rover recovers. Because
the fault acts on the hidden truth while the trust model observes only
outcomes, we can measure a \emph{recovery time}: how many cycles elapse
between the end of the fault and the point at which the learned trust
value and the assignment behaviour reflect the recovered reality. This
metric captures the responsiveness of the closed loop, not merely its
steady-state accuracy.

Complementing faults, an \emph{edge lifecycle tracker} records when links
are removed and re-added by the pruner across cycles. A re-added edge
receives a probation boost: it is provisionally trusted more than its
raw features would suggest, so that the graph can recover quickly from a
pruning decision that later turns out to have been too aggressive. The
tracker also yields the topology churn statistics reported alongside the
performance metrics, letting us distinguish a pruner that is stable but
wrong from one that is adaptive but noisy.

\subsection{Cloud Amortization and Staleness}

Solving the constrained QUBO is the most expensive stage of the pipeline,
and on a real deployment it would run on a cloud tier rather than on the
rovers. The harness models this with a \emph{solve interval} $k$: the
QUBO is solved every $k$ cycles, and between solves the most recent kept
set is reused verbatim. This creates a controlled notion of staleness.
With $k = 1$ the pruner always acts on fresh features; with larger $k$
the fleet operates on a pruning decision that may no longer match the
observed world. A dedicated staleness mode sweeps the interval and
quantifies the margin the pipeline retains, which is the quantity a
deployment engineer actually needs when sizing the cloud link budget.

\subsection{Metrics}

The harness reports a fixed battery of metrics for every run:

\begin{itemize}
  \item \textbf{Task success rate}, the primary mission metric, computed
    under the causal delivery model.
  \item \textbf{Per-stage wall-clock latency}, measured with a
    median-of-repeats timer. We deliberately avoid coarse clocks: the
    heuristic baselines run in microseconds, and reporting them as zero
    would conceal exactly the cost asymmetry between the QUBO pruner and
    its competitors that any honest comparison must show.
  \item \textbf{Trust-versus-truth Spearman correlation}, which grades
    the trust model against the hidden reliability field without
    requiring the trust values to be calibrated probabilities.
  \item \textbf{Topology churn}, the rate at which the kept graph
    changes across cycles.
  \item \textbf{Edge lifecycle counts}: additions, removals, and
    re-additions under probation.
  \item \textbf{Fault recovery time}, as defined above.
\end{itemize}

\subsection{Statistical Protocol}
\label{sec:sim-stats}

Every comparison claim in this report carries two statistical
instruments, computed on paired per-seed differences so that common
randomness across arms is cancelled rather than mistaken for signal.
First, we report a \emph{paired bootstrap confidence interval} with
$10\,000$ resamples, which quantifies the size and precision of the
difference. Second, we report a \emph{permutation test} with $20\,000$
permutations, which yields an exact finite-sample $p$-value under the
null hypothesis that the two arms are exchangeable. A difference is
described as significant only when the interval excludes zero and the
permutation test agrees; differences that fail this test, including some
that would flatter our own method, are reported as not significant.

\subsection{Run Modes}

The harness exposes five run modes, each corresponding to a class of
experiment reported in later chapters:

\begin{description}
  \item[Single run.] One mission with a per-cycle log, used for
    inspection, debugging, and the worked traces in this report.
  \item[Benchmark.] A pruner comparison at a fixed operating point: 10
    rovers, 20 steps, and 10 seeds, with real per-stage timing for every
    pruner. The four pruners compared are summarised in
    Table~\ref{tab:pruners}.
  \item[Ablation.] Four assignment arms over the same harness: the full
    GNN pipeline, a no-graph arm that scores nodes without message
    passing, a trust-only arm, and a random arm. This isolates the
    contribution of each modelling choice.
  \item[Staleness.] A sweep of the solve interval described above,
    measuring the margin retained under cloud amortization.
  \item[Training.] A held-out evaluation of the trained readout against
    the fixed-weight readout, so that the training result (which is
    negative) is measured on data the fit never saw.
\end{description}

\begin{table}[htbp]
  \centering
  \caption{The four pruning strategies compared in benchmark mode.}
  \label{tab:pruners}
  \begin{tabular}{@{}lp{8.2cm}@{}}
    \toprule
    Pruner & Selection rule \\
    \midrule
    QUBO (production) & Constrained optimization over utility, degree,
      budget, and connectivity, solved by projected annealing \\
    Full graph & Keep every candidate edge; no pruning \\
    Quality threshold & Keep edges with link quality at least $0.4$ \\
    Top-$K$ per node & Keep the best fraction $0.5$ of each node's edges \\
    \bottomrule
  \end{tabular}
\end{table}

Taken together, the causal delivery model, the hidden spatially
correlated ground truth, the fault and staleness machinery, and the
paired statistical protocol make the harness an audit instrument rather
than a demonstration stage. It is designed so that the pipeline can
lose, and the chapters that follow report the cases where it does with
the same prominence as the cases where it wins.
